%% Special Preface (1998) -- original pages v-vi
\fmhead[5]{SPECIAL PREFACE}{Special Preface}{specpref}{SPECIAL PREFACE}

\begingroup\small
The present edition differs from the original one (published in 1968) by:
\begin{itemize}
\item the inclusion of short notes giving references to new results;
\item a supplementary bibliography.
\end{itemize}
Otherwise, the text has been left unchanged, except for the correction of a
few misprints.

The added bibliography does not claim to be complete. Its aim is just to
help the reader get acquainted with some of the many developments of the
past twenty years (for those prior to 1977, see also the survey
\cite{serre-kyoto}). Among these developments, one may especially mention
the following:

\subsection*{\normalfont\bfseries \textit{l}-adic representations associated to abelian varieties over number fields}

Deligne (cf.~\cite{deligne-hodge}) has proved that Hodge cohomology
classes behave under the action of the Galois group as if they were
algebraic, thus providing a very useful substitute for the still unproved
Hodge conjecture.

Faltings (\cite{faltings-endl}, see also Szpiro \cite{szpiro} and
Faltings--Wüstholz \cite{faltings-wustholz}), has proved Tate's conjecture
that the map
\[
\Hom_K(A,B)\tens\Z_l \longrightarrow \Hom_{G_l}(T_l(A),\,T_l(B))
\]
is an isomorphism ($A$ and $B$ being abelian varieties over a number field
$K$), together with the semi-simplicity of the Galois module
$\Q_l\tens T_l(A)$ and similar results for $T_l(A)/lT_l(A)$ when $l$ is
large enough. These results may be used to study the structure of the
Galois group of the division points of $A$, cf.~\cite{serre-resume85}. For
instance, if $\dim A$ is odd and $\End A=\Z$, one can show that this
Galois group has finite index in the group of symplectic similitudes; for
elliptic curves, i.e. $\dim A=1$, this was already proved in
\cite{serre-proprietes}.

\subsection*{\normalfont\bfseries Modular forms and \textit{l}-adic representations}

The existence of \textit{l}-adic representations attached to modular
forms, conjectured in the first edition of this book, has been proved by
Deligne (\cite{deligne-formes}, see also Langlands \cite{langlands-modular}
and Carayol \cite{carayol}). This has many applications for instance to the
Ramanujan conjecture (Deligne) and to congruence properties (Ribet
\cite{ribet-ladic}, \cite{ribet-nebentypus}; Swinnerton-Dyer \cite{swd};
\cite{serre-ramanujan}, \cite{serre-congruences}). Some generalizations are
known (e.g. Carayol \cite{carayol}; Ohta \cite{ohta}; Wiles \cite{wiles}),
but one can hope for much more, in the setting of ``Langlands' program'':
there should exist a diagram
\[
\begin{array}{c@{\hspace{3em}}c@{\hspace{3em}}c}
\text{motives} & & \text{automorphic representations of}\\
\big\vert & \longrightarrow & \text{reductive groups}\\[0.15\baselineskip]
\text{rational \textit{l}-adic representations} &
\end{array}
\]
where the vertical line is (essentially) bijective and the horizontal arrow
injective with a precise description of its image (Deligne
\cite{deligne-valeurs}; Langlands \cite{langlands-automorphic};
\cite{serre-kyoto}). Such a diagram would incorporate, among other things,
the conjectures of Artin (on the holomorphy of L-functions) and
Taniyama--Weil (on elliptic curves over $\Q$). Chapters II and III of the
present book, supplemented by the results of Deligne (\cite{deligne-motifs})
and Waldschmidt (\cite{henniart}, \cite{waldschmidt}), may be viewed as
a partial realization of this ambitious program in the abelian case.

\subsection*{\normalfont\bfseries Local theory of \textit{l}-adic representations}

Here the ground field $K$, instead of being a number field, is a local
field of residue characteristic $p$. The most interesting case is
$\mathrm{char}\,K=0$ and $p=l$, especially when a Hodge--Tate decomposition
exists: indeed this gives precious information on the image of the inertia
group (Sen \cite{sen}; \cite{serre-hodgetate}; Wintenberger
\cite{wintenberger}). When the \textit{l}-adic representation comes from a
divisible group or an abelian variety, the existence of such a
decomposition is well known (Tate \cite{tate-pdiv}; see also Fontaine
\cite{fontaine-tatemod}); for representations coming from higher dimension
cohomology, it has been proved recently by Fontaine--Messing (under some
restrictions, cf.~\cite{fontaine-messing}) and Faltings
(\cite{faltings-hodge}). The results of Fontaine--Messing are parts of a
vast program by Fontaine, relating Galois representations and modules of
Dieudonné type (over some ``Barsotti--Tate rings,'' cf.~\cite{fontaine-galmod},
\cite{fontaine-barsotti}, \cite{fontaine-repp}).

\medskip
\hfill \textsc{Jean-Pierre Serre}

\medskip
\hfill \textit{Paris, Spring 1988}
